\documentclass[10pt,twocolumn]{article}

% ------------------ PAQUETES ------------------
\usepackage[spanish]{babel}
\usepackage[utf8]{inputenc}
\usepackage[T1]{fontenc}
\usepackage{geometry}
\usepackage{graphicx}
\usepackage{amsmath, amssymb}
\usepackage{siunitx}
\usepackage{authblk}
\usepackage{abstract}
\usepackage{caption}
\usepackage{titlesec}
\usepackage{multicol}
\usepackage{hyperref}
\usepackage{currfile}
\usepackage{float}

\geometry{margin=1.8cm}

% ------------------ FORMATO DE SECCIONES ------------------
\titleformat{\section}{\large\bfseries}{\thesection}{1em}{}
\titleformat{\subsection}{\normalsize\bfseries}{\thesubsection}{1em}{}

% ------------------ TÍTULO ------------------
\title{\textbf{Estrellas variables y Machine Learning}\\Trabajo con métodos de machine learning para la correcta identificación de estrellas variables\\Proyecto final del curso}


% ------------------ AUTORES ------------------
\author[1]{Bruno Abello Laurent}

\affil[1]{Universidad de Los Andes, Departamento de Física, Colombia}

% Correos (puedes modificar)
\renewcommand\Authands{, }
\date{}

\begin{document}

\twocolumn[
\maketitle
\begin{center}
\small
\textit{Correo institucional:} \href{mailto:{b.abello@uniandes.edu.co}}{b.abello@uniandes.edu.co} \\
\textit{Repositorio de GitHub: } \href{https://github.com/Brunelleschk1/Personal-projects/tree/main/Proyectos\%20de\%20Astronom\%C3\%ADa/Laboratorio\%20Final}{Repositorio del laboratorio}
\end{center}

\begin{abstract}
\noindent

\end{abstract}

\noindent\textbf{Palabras clave:}
\vspace{0.5cm}
]

% ------------------ CONTENIDO ------------------

\section{Introducción}

La astrofísica observacional moderna es probablemente una de las ramas de la física que menos se parece a lo que era hace cien años.
Hoy en día, la mayor parte del tiempo de un astrofísico corriente transcurre frente a una pantalla con un proyecto de Machine learning 
en curso. Aunque sea un método extremadamente popular actualmente, lo cierto es que esta es una de las herramientas más confiables del 
físico, ya que dada la lógica correcta, esta es la manera más rápida y más efectiva de tratar los miles de datos que capturan los telescopios
modernos, y que se hacen públicos gracias al formato informático. 

Este es el último en una serie de once informes planeados sobre la astrofísica. La idea de esta serialización es servir de aproximación 
a todo lo que hace un astrofísico moderno, desde las implementaciones de cómputo más básicas hasta los ejercicios de pensamiento más 
complejos que se tienen en el dominio. Después de pasar unas semanas hablando de cúmulos globulares, y utilizando un poco de machine learning
para buscar cúmulos nuevos en la vecindad del Sol, pasamos a trabajar las estrellas variables: un análisis básico de algunos tipos de estrellas 
variables fue uno de los once trabajos realizados. 

En este orden de ideas, este proyecto final trata las estrellas variables desde una perspectiva que no se había planteado en trabajos anteriores: 
¿Y si no sabemos qué tipo de estrella tenemos en los datos? A esto se enfrentan los astrofísicos muy a menudo, puesto que un telescopio no puede 
decirnos qué es una estrella, sólo cómo se comporta. El trabajo de encontrar qué tipo de estrella tenemos es algo tedioso, en teoría. Ver su curva 
de luz, su comportamiento, pensar en si encontramos su periodo correcto o si se trata solamente de un error de la máquina, son cosas que nunca 
vamos a poder solucionar del todo. 

Lo que sí podemos hacer es enseñarle a una máquina cómo se comportan las estrellas reales y decirle que 
clasifique una serie de datos por nosotros. Este es precisamente nuestro trabajo, entre estrellas eclipsantes, pulsantes, de largo periodo y estrellas
no clasificadas. Gracias al algoritmo de Random Forests, y a unas 6000 estrellas
repartidas entre set de entrenamiento y de clasificación, este trabajo presenta ventajas y desventajas del machine learning aplicado a la 
astrofísica observacional.

\subsection{Objetivos}

Los objetivos propuestos para este informe son los siguientes:
\begin{itemize}
    \item Aprender a implementar un método de machine learning con datos etiquetados.
    \item Aprender a leer los resultados de dicho método.
    \item Aprender a interpretar los resultados con el objetivo claro de mejorar la calidad de la información que trabaja el método.
    \item Aprender a pensar como astrofísico antes de implementar alguna función o método computacional para el análisis de datos.
    \item Entender por qué el conocimiento adquirido como físico es tan importante para el uso del machine learning y demás métodos computacionales.
\end{itemize}

\section{Procedimiento}

Este trabajo tuvo un proceder distinto a los anteriores. De las tres semanas que se trabajó, solo en los últimos días se hizo algo de 
machine learning. Las primeras dos semanas fueron de revisión de datos y de ideación de features para el método. Hacer un trabajo de estos 
no es una tarea sencilla, ya que se requiere de un gran conocimiento para entender qué es una buena feature independiente para un 
dataset específico. En este caso, todos los estudiantes empezamos con periodo, que luego se cambió por logaritmo de periodo, que parece
separar mejor las estrellas, y luego tomamos color V-I. Adicionalmente, muchos optamos por utilizar skewness, ya que presenta 
mejores resultados que la kurtosis, y a modo personal se escogió la diferencia de fase entre un máximo y el siguiente mínimo. 

Cada una de estas features se estudió cuidadosamente y se realizaron extensivas gráficas y comparaciones para determinar si 
eran features dignas de ser usadas. Esto sucedió principalmente con la feature personal, ya que la idea provino de que algunas
estrellas variables tienen una curva de luz en dientes de sierra, y otras son más sinusoidales: si se pudiera capturar esta diferencia
mediante el offset de la curva, tendríamos una feature potente. Sin embargo, se terminó tomando únicamente la separación entre un 
máximo y su siguiente mínimo ya que el contrario, y la diferencia entre mínimos para detectar estrellas binarias, no mostraron separación
significativa entre las clases de estrellas.

Se realizaron gráficos 3D para comparar separaciones de estrellas según features, y se intentaron varias features que no se incluyeron en 
el método final. 

\section{Problemas presentados y soluciones aplicadas}

\begin{enumerate}
    \item Entender por qué algunas features son mejores que otras resultó bastante complicado, especialmente entender por qué features que parecen tan evidentes en una gráfica no se 
    traducen a algo en el código.
    \item Como se dijo, algunas features que parecían prometedoras no fueron utilizadas porque, a la hora de la verdad, no representaban en números lo que se ve en un gráfico. 
    \item La feature de periodo presentó problemas con las estrellas eclipsantes, que tienen dos picos de distinta intensidad, y con las LPVs, que tienen un rango enorme de periodos. 
    No hay una solución clara al respecto, puesto que para corregir el error hay que saber que se trata de una estrella de este tipo, así que se esperan más confusiones
    en estas categorías que en las demás. 
    \item Las iteraciones e intentos en el código lo volvieron bastante lento. Se puede optimizar, pero no es necesario por ahora, puesto que, dentro de todo, random forest es de las 
    partes más rápidas del código. 
\end{enumerate}


\section{Resultados}

Puesto que este es un trabajo extenso, esta sección se divide en la parte del trabajo sobre features y el trabajo sobre los resultados.

\subsection{Lomb-Scargle}

(Fig: ~\ref{STAR29OGLE}) 

\begin{figure}[h]
    \centering
    \includegraphics[width=1\linewidth]{../OGLE_Results/Star_29.png}
    \caption{Estrella 29 de las analizadas en OGLE. Se nota ruido alrededor de la frecuencia real de la estrella, y se ven 
    caídas en la curva de luz que no parecen ser consistentes sobre todo el periodo de observación de la estrella.}
    \label{STAR29OGLE}
\end{figure}

\section{Conclusiones}


\section*{Comentarios al margen}


\end{document}